% =============================================================================
% practica4.tex --- Chapter: Practice 4
%
% Compile standalone:
%   pdflatex -shell-escape practica4.tex   (run twice for cross-references)
%
% Compile as part of the full report:
%   pdflatex -shell-escape __memoria.tex   (from docLaTex root)
% =============================================================================

\documentclass[../../__memoria.tex]{subfiles}

\begin{document}
\standalonecover
\pagestyle{cuerpo}

\graphicspath{{\subfix{../../Imagenes/}}{../../Imagenes/}{\subfix{img/}}{img/}}

\chapter[Práctica 4: Clasificación y Matriz de Costes]{Cuaderno de Laboratorio --- Práctica 4: Clasificación y el Valor del Error (Matriz de Costes)}

\section*{Objetivo}
Entrenar un clasificador probabilístico para un problema de detección de fraude bancario y estudiar cómo el umbral de decisión afecta a los errores cometidos y al coste total. El objetivo central no es únicamente medir el rendimiento predictivo, sino comprender el impacto económico real de las decisiones del modelo a partir de una \textbf{matriz de costes asimétrica}.

\section*{Descripción del problema}
Se trabaja con un problema de clasificación binaria sobre el dataset \texttt{fraude\_transacciones.csv}:
\begin{itemize}
  \item \textbf{Clase 0}: transacción normal.
  \item \textbf{Clase 1}: transacción fraudulenta.
\end{itemize}

Los errores tienen costes muy distintos:
\begin{itemize}
  \item \textbf{Falso Negativo (FN)}: no detectar un fraude $\rightarrow 200\,\text{€}$
  \item \textbf{Falso Positivo (FP)}: bloquear una transacción legítima $\rightarrow 10\,\text{€}$
\end{itemize}

La decisión final no debe evaluarse solo por el número de aciertos, sino por el \textbf{coste económico} asociado a los errores.

% =============================================================================
\section[Tarea 1: Umbral por defecto]{Tarea 1 --- Análisis con el umbral por defecto (\texorpdfstring{$\tau = 0.5$}{tau=0.5})}
% =============================================================================

\subsection{Modelo utilizado}
Se ha entrenado una \textbf{Regresión Logística} (\texttt{LogisticRegression}, \texttt{max\_iter=1000}) sobre el 70\% del dataset (partición estratificada, \texttt{random\_state=42}) y evaluado sobre el 30\% restante.

\subsection{Matriz de confusión}

\begin{figure}[H]
  \centering
  \includegraphics[width=0.52\textwidth]{matriz_confusion_umbral_0_5}
  \caption{Matriz de confusión para el umbral estándar $\tau = 0.5$.}
  \label{fig:mc_05}
\end{figure}

A partir de la matriz de confusión detallada en la Figura~\ref{fig:mc_05}, se obtienen los siguientes valores:
\begin{itemize}
  \item Verdaderos Positivos (TP): \textbf{16}
  \item Verdaderos Negativos (TN): \textbf{501}
  \item Falsos Positivos (FP): \textbf{1}
  \item Falsos Negativos (FN): \textbf{22}
\end{itemize}

\subsection{Cálculo del coste total}
\[
  \text{Coste Total} = (FP \cdot 10) + (FN \cdot 200) = (1 \cdot 10) + (22 \cdot 200) = \mathbf{4410\,\text{€}}
\]

\subsection{Interpretación}
Con el umbral por defecto de 0.5, el modelo solo cometió 1 falso positivo pero dejó pasar 22 fraudes. El coste total fue de 4410\,€. Me llamó la atención que un umbral que parece razonable (0.5) pueda ser tan malo cuando los costes de los errores no son simétricos: cada fraude no detectado cuesta 200\,€, veinte veces más que un falso positivo (10\,€).

% =============================================================================
\section[Tarea 2: Probabilidades y Umbrales]{Tarea 2 --- Estudio de probabilidades y umbrales}
% =============================================================================

\subsection{Análisis de distintos umbrales}
Se han extraído las probabilidades de la clase fraude para cada observación y se han evaluado umbrales $\tau \in [0, 1]$ en pasos de 0.05. La Tabla~\ref{tab:umbrales_coste} resume los resultados para los primeros tramos y el extremo superior.

\begin{table}[H]
  \centering
  \caption{Resumen de FP, FN y Coste Total para distintos umbrales de decisión.}
  \label{tab:umbrales_coste}
  \begin{tabular}{cccc}
    \toprule
    \textbf{Umbral} & \textbf{FP} & \textbf{FN} & \textbf{Coste Total (€)} \\
    \midrule
    0.00            & 502         & 0           & 5020                     \\
    0.05            & 112         & 9           & \textbf{2920}            \\
    0.10            & 42          & 14          & 3220                     \\
    0.15            & 22          & 15          & 3220                     \\
    0.20            & 13          & 15          & 3130                     \\
    0.25            & 9           & 16          & 3290                     \\
    0.30            & 4           & 17          & 3440                     \\
    0.40            & 1           & 20          & 4010                     \\
    0.50            & 1           & 22          & 4410                     \\
    1.00            & 0           & 38          & 7600                     \\
    \bottomrule
  \end{tabular}
\end{table}

\subsection{Observaciones}
Existe un claro compromiso (\textit{trade-off}) entre los dos tipos de error al variar $\tau$:
\begin{itemize}
  \item \textbf{Cuando el umbral baja} (hacia 0): se exige menos confianza para clasificar como fraude. Caen los Falsos Negativos (detectamos más fraudes), pero aumentan drásticamente los Falsos Positivos (bloqueamos más transacciones legítimas).
  \item \textbf{Cuando el umbral sube} (hacia 1): las alertas son mucho más selectivas. Los Falsos Positivos desaparecen, pero los Falsos Negativos se disparan (dejamos escapar la mayoría de los fraudes).
\end{itemize}
El umbral de decisión debe ajustarse según la penalización económica de cada error, en lugar de asumir un valor estático por defecto.

{\sloppy
Lo interesante de trabajar con las probabilidades directamente es que puedes decidir tú dónde poner el límite según lo que quieras priorizar. Por ejemplo, una transacción con un 0.30 de probabilidad de fraude se clasificaría como normal con el umbral de 0.5, pero si bajas el umbral a 0.20 ya la detectas. La pregunta es cuántos falsos positivos estás dispuesto a aceptar a cambio de detectar más fraudes.

En este caso, como el coste de un falso negativo es 20 veces el de un falso positivo, parece razonable bajar el umbral bastante por debajo de 0.5. Eso es justo lo que se ve en la tabla: el coste mínimo se alcanza en $\tau = 0.05$, muy lejos del valor por defecto.
\par}


% =============================================================================
\section[Tarea 3: Optimización del coste]{Tarea 3 --- Optimización del coste}
% =============================================================================

\subsection{Gráfica de coste frente a umbral}

\begin{figure}[H]
  \centering
\includegraphics[width=0.78\textwidth]{coste_vs_umbral}
  \caption{Coste total (€) en función del umbral de decisión. Como se aprecia en la Figura~\ref{fig:coste_umbral}, el punto óptimo se alcanza en $\tau^* = 0.05$.}
  \label{fig:coste_umbral}
\end{figure}

\subsection{Umbral óptimo}
En base a la gráfica anterior, identificamos el punto de menor coste:
\[
  \tau^* = 0.05 \qquad \text{Coste mínimo} = 2920\,\text{€}
\]

\subsection{Interpretación}
El umbral óptimo ($\tau^* = 0.05$) es muy inferior a 0.5. Esto tiene todo el sentido dada la fuerte asimetría de la matriz de costes: el coste de un FN (200\,€) es 20 veces el de un FP (10\,€). Para minimizar el coste total conviene ser muy sensible a los fraudes, aunque eso suponga aceptar más falsos positivos. La estrategia más razonable para el banco es \textbf{conservadora}: es preferible molestar puntualmente a un cliente legítimo mediante un bloqueo temporal que permitir que un fraude se consume sin ser detectado.

% =============================================================================
\section[Tarea 4: Curvas ROC y PR]{Tarea 4 --- Curvas ROC y Precision-Recall}
% =============================================================================

\subsection{Curva ROC}

\begin{figure}[H]
  \centering
  \includegraphics[width=0.62\textwidth]{curva_roc}
  \caption{Curva ROC del clasificador. AUC-ROC $= 0.8442$.}
  \label{fig:roc}
\end{figure}

\subsection{Curva Precision-Recall}

\begin{figure}[H]
  \centering
  \includegraphics[width=0.62\textwidth]{curva_precision_recall}
  \caption{Curva Precision-Recall del clasificador. Average Precision $= 0.6551$.}
  \label{fig:pr}
\end{figure}

\subsection{Discusión}
Al analizar un clasificador sobre un dataset desbalanceado, la curva Precision-Recall suele ser mucho más informativa que la ROC (Figura~\ref{fig:roc} frente a Figura~\ref{fig:pr}). El motivo es estructural: el \textit{False Positive Rate} de la curva ROC toma como denominador el total de negativos, que en este problema es muy numeroso. Como consecuencia, incrementos apreciables de falsos positivos apenas desplazan el FPR y el AUC-ROC se mantiene elevado (0.8442) pese a que el modelo comete errores relevantes sobre la clase fraude. La curva PR evita este efecto porque la Precisión (TP / (TP + FP)) cae en cuanto aumentan los FP, lo que la hace sensible exactamente a la calidad de recuperación de la clase minoritaria que interesa controlar.


% =============================================================================
\section{Decisión final}

Viendo los resultados, creo que lo más sensato es fijar $\tau = 0.05$ en lugar del umbral por defecto. El coste baja de 4410\,€ a 2920\,€, un ahorro del 34\,\%. Es cierto que los falsos positivos suben de 1 a 112, pero como cada falso negativo cuesta 20 veces más (200\,€ frente a 10\,€), compensa: se pasa de 22 fraudes no detectados a solo 9.

Eso sí, este umbral está calculado sobre los datos de prueba de esta práctica. Si fuese un sistema real habría que probarlo con datos nuevos para ver si se mantiene.

% =============================================================================
\section{Conclusión final}
% =============================================================================

Lo principal que me llevo de esta práctica es que el umbral por defecto de 0.5 no vale para todos los problemas. En este caso, con costes asimétricos, usarlo sin pensar suponía un coste de 4410\,€, cuando bajándolo a 0.05 se reducía a 2920\,€. La diferencia está en que un falso negativo cuesta 200\,€ y un falso positivo solo 10\,€, así que merece la pena arriesgarse a bloquear más transacciones legítimas con tal de no dejar escapar fraudes.

También me sirvió para ver que el AUC-ROC puede dar una visión demasiado optimista cuando las clases están desbalanceadas, mientras que la curva Precision-Recall muestra mejor cómo de bien está funcionando el modelo en la clase minoritaria. En mi opinión, ambas métricas son útiles pero ninguna debería mirarse sin tener en cuenta el coste real de cada tipo de error.

\end{document}
